\section{Self-dual stable sources reduce to odd rank two}
\label{sec:kt-selfdual-source-exclusion}

The determinant normalization in the preceding adjoint theorem is
not automatic for an arbitrary proposed two-dimensional origin.
The following automorphic descent first produces an odd origin for
every self-dual stable source.  The separate scalar-normalization
lemma then gives the exact ordinary rank-two exclusion.  In this
section, self-dual means $R^\vee\simeq R$ on $G_E$, with no Tate
twist; it is stronger than conjugate self-duality.

\begin{theorem}[KT-SELFDUAL-SOURCE-ODD-DESCENT]
Let $p>3$, let $E=\mathbf Q(\sqrt D)$ be imaginary quadratic with
$p\nmid D$, put $\eta=\chi_D$, and let $e\ge2$ be even with
$(p-1)\nmid2e$.
Let $R$ be a continuous absolutely irreducible three-dimensional
representation of $G_E$ over a finite extension of $\mathbf Q_p$.
Suppose that its residual semisimplification is
\[
                  1\oplus\omega^e\oplus\omega^{-e}.       \tag{SD.1}
\]
Suppose also that $R$ is attached, in the weight-zero unitary
normalization used in KT-U3-GENERAL-COVERAGE, to a regular algebraic
conjugate-self-dual cuspidal representation $\Pi$ of
$\mathrm{GL}_3(\mathbf A_E)$ whose complex infinity parameter is
\[
                 (z/\bar z)^e\oplus1\oplus(z/\bar z)^{-e}.
                                                               \tag{SD.2}
\]
If $R^\vee\simeq R$, there is a holomorphic cuspidal eigenform $f$
over $\mathbf Q$ of weight $e+1$, and a finite-coefficient
two-dimensional Galois representation $V_f$, such that
\[
              R\simeq\operatorname{Ad}^0(V_f)|_{G_E}.     \tag{SD.3}
\]
The representation $V_f$ is absolutely irreducible and is
projectively odd: $V_f(c)$ is not scalar for a complex conjugation
$c$.  No determinant normalization or fixed classical level is
asserted for $V_f$ at this stage.
\end{theorem}

\begin{proof}
\paragraph{The determinant is exactly one.}
Self-duality gives $(\det R)^2=1$.  Thus every value of $\det R$
is a sign.  Formula (SD.1) says that its reduction is one on every
element of $G_E$.  Reduction is injective on $\{\pm1\}$ because
$p>2$, so
\[
                              \det R=1.                  \tag{SD.4}
\]
This uses both actual self-duality and the residual determinant;
it does not infer (SD.4) from a residual statement alone.

The unitary-normalized attachment identifies the full unramified
Satake polynomials with the Frobenius polynomials of $R$ at almost
all places of $E$.  Self-duality and (SD.4), followed by strong
multiplicity one for $\mathrm{GL}_3$ and $\mathrm{GL}_1$, therefore
give
\[
                  \Pi^\vee\simeq\Pi,
                  \qquad\omega_\Pi=1.                   \tag{SD.5}
\]
The source's initial algebraic attachment is to
$R\varepsilon_p^{-1}$; the common twist in the already checked
unitary normalization is essential here.  No additional Tate
factor is inserted in (SD.5).

\paragraph{Cuspidal quadratic descent and its central character.}
Conjugate self-duality and (SD.5) imply
$\Pi^c\simeq\Pi$.  Arthur--Clozel, Chapter~3,
Theorem~4.2(d), printed p.~202, supplies a \emph{cuspidal}
representation $\Sigma$ of $\mathrm{GL}_3(\mathbf A_{\mathbf Q})$
with base change $\Pi$.  Thus no noncuspidal descent is being
silently used.  Theorem~5.1, pp.~212--213, makes this base change
strong at every finite and infinite place.

Base change of central characters and (SD.5) give
$\omega_\Sigma\circ N_{E/\mathbf Q}=1$.
Global class field theory for the quadratic extension gives
\[
                              \omega_\Sigma\in\{1,\eta\}.
                                                               \tag{SD.6}
\]
Both $\Sigma^\vee$ and $\Sigma$ base change to $\Pi$.
The precise fiber theorem, Arthur--Clozel Chapter~3,
Theorem~3.1, p.~201, or the fiber clause in Theorem~4.2(d),
gives $\Sigma^\vee\simeq\Sigma\otimes\eta^j$ with
$j\in\{0,1\}$.  If $j=1$, the central characters would satisfy
\[
                 \omega_\Sigma^{-1}
                          =\omega_\Sigma\eta^3,
\]
contradicting $\omega_\Sigma^2=1$ and $\eta\ne1$.
Consequently $\Sigma$ is self-dual.  Set
\[
                          \Sigma_0=\Sigma\otimes\omega_\Sigma.
                                                               \tag{SD.7}
\]
It remains cuspidal and self-dual, has central character
$\omega_\Sigma^4=1$, and still base changes to $\Pi$.

\paragraph{The real parameter gives a holomorphic rank-two origin.}
Strong base change at infinity identifies the restriction of the
real Weil parameter of $\Sigma_0$ to $\mathbf C^\times$ with
(SD.2).  In particular $\Sigma_0$ is regular algebraic, and its
three distinct exponents are exactly $(e,0,-e)$.

Ramakrishnan's \emph{An exercise concerning the selfdual cusp forms
on $\mathrm{GL}(3)$}, Theorem~A, gives
\[
                            \Sigma_0\simeq\operatorname{Ad}(\pi)
                                                               \tag{SD.8}
\]
for a non-dihedral cuspidal representation $\pi$ of
$\mathrm{GL}_2(\mathbf A_{\mathbf Q})$; its central character
can be chosen of finite order.  The quadratic factor in that
theorem is the central character of $\Sigma_0$, so it is one.
Corollary~B gives regular algebraicity after the usual common
norm normalization.  The explicit archimedean compatibility
in Proposition~C(b) and formula (6b) also fixes the weight here:
the two real discrete-series characters have ratio
$(z/\bar z)^e$, so their discrete-series weight is $e+1$.
Since $e\ge2$, this is a genuine discrete series.  Thus $\pi$
is represented by a holomorphic cuspidal eigenform $f$ of that
weight (choosing its holomorphic vector), in the ordinary
unitary normalization of a classical form.

Let $V_f$ be its usual finite-coefficient Galois representation.
If $\xi$ is the finite nebentypus, then
\[
                 \det V_f=\xi\varepsilon_p^e,
                 \qquad \xi(-1)=(-1)^{e+1}=-1.           \tag{SD.9}
\]
Hence $\det V_f(c)=-1$, whereas $V_f(c)^2=1$; its two
eigenvalues are $1,-1$.  In particular it is projectively odd.

At almost every place of $E$, strong base change, (SD.8),
Gelbart--Jacquet's adjoint local compatibility, and the classical
Galois attachment give the same full characteristic polynomial
for $R$ and $\operatorname{Ad}^0(V_f)|_{G_E}$.  Chebotarev and
Brauer--Nesbitt identify their semisimplifications.  Since $R$
is irreducible of dimension three, the other representation
is itself irreducible and isomorphic to $R$, giving (SD.3).
It also follows that $V_f$ is absolutely irreducible: a
reducible two-dimensional representation has reducible adjoint.
Thus no separate all-primes big-image assertion is needed.
\end{proof}

\begin{corollary}[KT-SELFDUAL-SOURCE-UNIT-EXCLUSION]
Assume the base hypotheses on $p,E,e,R,\Pi$ and (SD.1)--(SD.2)
in KT-SELFDUAL-SOURCE-ODD-DESCENT, without assuming that $R$ is
self-dual.  Suppose in addition that $R$ is unramified at all finite places of $E$
away from $p$, and at a place above $p$ it has a stable full flag
with ordered inertial characters
$(\varepsilon_p^e,1,\varepsilon_p^{-e})$.
The unit condition
\[
                          B_{1,\eta\omega^e}\in\mathbf Z_p^\times
                                                               \tag{SD.10}
\]
implies $R^\vee\not\simeq R$.
\end{corollary}

\begin{proof}
If $R$ were self-dual, (SD.3) would give the projectively odd
origin $V_f$.  Its adjoint has precisely the residual,
away-prime and ordered ordinary properties in the hypotheses.
KT-ODD-ADJOINT-SCALAR-NORMALIZATION supplies a global scalar
twist $V'$ with the residual pair $1\oplus\eta\omega^e$,
an unramified ordinary quotient at $p$, and unramified
restriction to $G_E$ away from $p$.
KT-ORDINARY-RANK2-UNIT-EXCLUSION, with $a=e$, contradicts
(SD.10).
\end{proof}

\paragraph{Exact actual-source application and remaining gate.}
For the actual $N=2$, $\psi=1$ source and GD, under the carried
closed source convention KT-ONECHAR-CLOSED, the Type~I clause
in \texttt{FOneCharU.lean}, lines 47--54, supplies
\texttt{Adm3U\_s}.  Its base predicate
\texttt{AfStmt.lean}, lines 29--32, supplies irreducibility,
the algebraic automorphic attachment and the local level.
For the original ordered flag, the exact definitions are
\texttt{AutomorphicForms/BlochKatoSelmerVanishes.lean},
\texttt{IsFlagOn} at lines 37--39, \texttt{KillsMu} at
lines 66--67, \texttt{LevelAt} at lines 102--107, and
\texttt{Level.above} at lines 111--114.
Every character satisfying \texttt{KillsMu 2} is trivial because
$\mu_2\subset\mathbf Q$.  Hence that level clause is an actual
decomposition-group stable upper flag with exactly the ordered
inertial characters $(\varepsilon_p^e,1,\varepsilon_p^{-e})$;
it is not inferred just from potential crystallinity.
\texttt{IsAutomorphicOverImagQuad.lean}, lines 35--36, is
explicitly conjugate-self-dual and cuspidal.
The checked common twist and \texttt{KeepLevel.lean},
lines 30--82 and 95--101, give the unitary-normalized attachment
and (SD.2).  The actual residual matrix from
\texttt{IsPtOf.lean}, lines 124--166, gives (SD.1), and
\texttt{PtInertiaAway.lean}, lines 54--69, gives the required
away-prime unramifiedness.  Finite coefficients follow from
\texttt{TanDet.lean}, lines 212--216, exactly as in
KT-AUXILIARY-UNIT-NO-TYPEII.

Accordingly, the exponent-$e$ auxiliary unit excludes
\emph{every self-dual actual Type~I source}, including any
adjoint origin over $E$ because an adjoint is self-dual.
The exponent-$2e$ unit already excludes every actual Type~II
source by KT-AUXILIARY-UNIT-NO-TYPEII.
Under these two units, any remaining actual source is a
non-self-dual stable rank-three source.  Conjugate self-duality
then says that its conjugate is its distinct dual.  This is a
strictly narrower source class; no assertion excludes that
remaining class or proves the original strict Selmer vanishing.

\paragraph{Primary source and review boundary.}
The new published input is D.~Ramakrishnan, \emph{Indian Journal
of Pure and Applied Mathematics} \textbf{45} (2014), 777--785,
\url{https://doi.org/10.1007/s13226-014-0088-1}.
The whole five-page author version at
\url{https://math.uchicago.edu/~fcale/papers/Ramakrishnan.pdf}
was read: Theorem~A and Corollary~B on author pp.~1--2,
Proposition~C(b) on p.~3, formula (6b) on p.~4, and the
complete Corollary~B proof on pp.~4--5.  The foundational
descent and transfer theorems it cites are imported through
Ramakrishnan's theorem, not freshly reconstructed here.
The author-version and published pagination are deliberately
distinguished.  Arthur--Clozel's precise quadratic descent,
fiber and strong-base-change statements are the same checked
primary inputs used earlier, reread here for the descent scope.
The scalar-normalization lemma and the descent and exclusion
proofs above have completed whole-proof and exact-source review.
No new Lean kernel build is claimed.
